\subsection{A unified description}

The three studies support one description of prediction-market prices. A binary contract's fair
price is a digital option on a latent state, written either as a score, as a numeric outcome, or as
the waiting time to an event. The state is revealed according to an information clock, with
short-lived random variation in the rate of revelation and a large share of discrete jumps:
\begin{align*}
  \text{single contract or ladder leg:}\quad &\pr_t=\Prob(Y_T>K\mid\F_t),\\
  &\dd Y_t=\alpha_t\,s(t)\,\dd W_t+J_t\,\dd N_t,\qquad \dd\alpha_t=\nu\alpha_t\,\dd B_t,\\
  \text{deadline contract:}\quad &\pr_t=1-\E_t\Big[e^{-\int_t^T h_u\dd u}\Big],\qquad
  h_u=h^{\det}(T-u)\,\xi_u ,
\end{align*}
where $h^{\det}$ is a deterministic hazard profile that rises into the deadline and $\xi$ a
log-diffusion. \Cref{tab:disc-facts} collects the estimates that pin down each component.

\begin{table}[t]
\centering
\begin{threeparttable}
\caption{Stylized facts by model component and contract type.}\label{tab:disc-facts}
\small
\begin{tabular}{@{}p{0.2\textwidth}p{0.36\textwidth}p{0.36\textwidth}@{}}
\toprule
Component & Estimate & Where it comes from \\
\midrule
Variance budget & realized variance over $\pr_0(1-\pr_0)$: 0.93--1.01 & weather, sports (\cref{sec:res-dyn}); 1.19 for deadlines (\cref{sec:res-haz}) \\
Clock & uniform in-game; back-loaded weather; flat pre-game; hazard rising into deadline & \cref{sec:res-dyn,sec:res-lad,sec:res-haz} \\
Backbone & $\beta\approx0.3$--$0.8$; logit-normal rejected & \cref{sec:res-dyn} \\
Stochastic volatility & half-life 7--21\,h (weather about 40\,h), $\nu\approx0.3$--$0.5$ per $\sqrt{\text{day}}$ & \cref{sec:res-dyn} \\
Jumps & 55--85\% of hourly variance; about half of in-game variance at 5--15\,s & \cref{sec:res-dyn,sec:res-hf} \\
Smile & convex, sticky in moneyness; tails too heavy & \cref{sec:res-lad} \\
Premia & implied tail mass, hazard and variance exceed realized & \cref{sec:res-lad,sec:res-haz} \\
\bottomrule
\end{tabular}
\end{threeparttable}
\end{table}

\subsection{Why clocks differ}

The information clock is the object that varies most across contract types, and it has a clear
economic reading: it is the schedule on which the world reveals the answer.
\begin{itemize}
  \item \emph{In-game sports.} Scoring opportunities arrive roughly uniformly in game time, so a
  uniform clock fits, exactly as in \citet{stern1994}.
  \item \emph{Pre-game sports.} Before the start almost nothing is learned, apart from occasional
  team news, so the clock runs at a few percent of the in-game rate.
  \item \emph{Weather.} Daily temperature extremes are revealed by observations late in the local
  day, giving a back-loaded clock with a pronounced time-of-day profile.
  \item \emph{Deadline contracts.} The implied hazard rises into the deadline. Part of this is real,
  since announcements do bunch late. But realized event timing is less concentrated than prices
  imply, so part of the rise is a premium (\cref{sec:res-haz}).
\end{itemize}

For pricing, the clock is to a binary contract what the term structure of implied variance is to
an option. A single contract's price carries no volatility information. A ladder's at-the-money
term structure and a deadline strip's forward hazards do, and both can be compared with the clock
measured from price paths.

\subsection{Jumps and the limits of the Gaussian model}

M0 gets the clock right and the shape wrong. Outside in-game sports, hourly standardised
innovations have excess kurtosis from tens to hundreds. Even within games, roughly half of
second-scale variance arrives in discrete moves. Several consequences follow.
\begin{itemize}
  \item \emph{Variance forecasts.} QLIKE and MSE rank the forecasting models differently: filtering
  the volatility state helps in quiet periods and hurts across jumps.
  \item \emph{Drift.} The outward drift of the probit price per unit of total variance falls short
  of the Gaussian prediction, as \cref{sec:model-m2} anticipates.
  \item \emph{Derivative pricing.} Any instrument that depends on the path, such as the
  variance and first-passage claims proposed by \citet{dalen2025}, must model jumps explicitly.
\end{itemize}

The basketball result is a useful benchmark. In the NBA, jump sizes scale with the Stern factor
$\varphi(\Zt)/\sqrt{\tauexp}$ almost exactly. In lower-scoring sports they scale less steeply.
This is consistent with scores that are rarer, larger and more state-dependent, for which
\citet{fengpolsonxu2016}'s Skellam model is the natural replacement.

\subsection{Backbone and smile dynamics}

The estimated backbone exponent of 0.3--0.8 places price volatility closer to the Wright--Fisher
(Jacobi) diffusion than to the logit-normal model that is common in practice. This independently
supports the Wright--Fisher component of \citet{xi2026}.

In ladders, the smile moves with the implied median: volatility at a fixed strike shifts as the
median moves, while at-the-money volatility does not slide along a fixed smile. In
\citet{derman1999}'s terms, the market treats the outcome distribution as a location--scale family
that is re-centred on news. That is natural when news shifts the expected outcome (a score, a
forecast update) without changing the remaining uncertainty much.

\subsection{A tail-and-time premium}

The one systematic departure from fair pricing appears in three forms:
\begin{itemize}
  \item implied tails across strikes are heavier than realized tails;
  \item implied hazards exceed realized hazards at every lead before a deadline, and prices decay
  more slowly than the hazard model requires;
  \item before and during games, ladders imply more variance than is subsequently realized.
\end{itemize}
In each case the market overprices the probability of an outcome far from the centre, or of an
event within a short window. This is the favourite--longshot bias of \citet{snowbergwolfers2010}
and \citet{burgi2025}, seen across strikes and across time.

Three features of the venue limit how far it can be corrected.
\begin{itemize}
  \item \emph{Tick grid.} The one-cent grid caps the gain from selling an overpriced 1--2\cents
  contract at about a cent.
  \item \emph{Fees.} Per-order rounding of fees makes small orders near the edges expensive.
  \item \emph{Collateral.} Selling a longshot requires posting nearly the full contract value as
  collateral.
\end{itemize}
In line with these limits, simple rules that sell the premium earn at most a few cents per contract
out of sample, and none is robust across samples. The premium is a stable feature of prices, not
an exploitable inefficiency at retail scale.

Two departures from the martingale property remain.
\begin{itemize}
  \item \emph{Weather.} Ladder medians for daily highs drift down overnight and up during the day.
  This may reflect a forecast input that under-predicts daytime maxima until observations arrive.
  \item \emph{Deadline contracts.} Prices fall faster than the martingale allows in their final
  hours, as stale prices catch up with the deadline.
\end{itemize}
Both are localised and small.

\subsection{Limitations}

The sample covers February to October 2026 for quotes and November 2026 for order books, a period
of rapid growth on the venue. Hourly candles hide intra-hour dynamics, and the base quote panel is
weighted toward actively traded markets. The order-book study covers two evenings and 50 game
markets.

We do not observe the game clock or the score. In-game time to expiry is measured to scheduled
settlement, which probably explains part of the cross-league dispersion. Event times of deadline
contracts are upper bounds inferred from prices and closing times.

Finally, we report many cells across categories, horizons and leads. We emphasise patterns that
hold in sign across the estimation and holdout samples, but individual cells should be read with
multiple testing in mind.
